% Gaussian elimination with partial pivoting (the pivoting steps are in red),
% innermost for loop collapsed into a whole-row operation.  Cropped to its own
% bounding box so it can be dropped onto a Xournal++ page as an image.
%
%   latexmk -pdf gauss-pivot-pseudocode.tex
\documentclass[border=6pt,varwidth=10cm]{standalone}

\usepackage{algorithm2e}
\usepackage{amsmath}
\usepackage{xcolor}

\RestyleAlgo{plain}          % no float, no ruled box
\DontPrintSemicolon
\SetKwComment{Comment}{// }{}
\SetKwInOut{Input}{input}
\SetKwInOut{Output}{output}

\begin{document}
\begin{minipage}{\linewidth}
\begin{algorithm}[H]
  \Input{$n \times n$ matrix $[A]$, right-hand side $\{b\}$}
  \Output{solution $\{x\}$ of $[A]\{x\} = \{b\}$}
  \BlankLine
  \tcp{forward elimination}
  \For{$c \leftarrow 1$ \KwTo $n-1$}{
    \textcolor{red}{\tcp{find the largest entry in column $c$}}
    \textcolor{red}{$p \leftarrow \operatorname{argmax}(\operatorname{abs}(A[c\!:\!n,\, c])) + c - 1$}\;
    \textcolor{red}{\tcp{swap that row into the pivot position}}
    \textcolor{red}{$A[c,\,:] \leftrightarrow A[p,\,:]$}\;
    \textcolor{red}{$b[c] \leftrightarrow b[p]$}\;
    \For{$r \leftarrow c+1$ \KwTo $n$}{
      $f \leftarrow A[r,c] \,/\, A[c,c]$\;
      $A[r,c\!:\!n] \leftarrow A[r,c\!:\!n] - f \cdot A[c,c\!:\!n]$\;
      $b[r] \leftarrow b[r] - f \cdot b[c]$\;
    }
  }
  \BlankLine
  \tcp{back substitution}
  \For{$r \leftarrow n, n-1, \ldots, 1$}{
    $s \leftarrow b[r]$\;
    \For{$j \leftarrow r+1$ \KwTo $n$}{
      $s \leftarrow s - A[r,j] \cdot x[j]$\;
    }
    $x[r] \leftarrow s \,/\, A[r,r]$\;
  }
\end{algorithm}
\end{minipage}
\end{document}
